\documentclass[11pt]{article}

\usepackage[margin=1in]{geometry}
\usepackage{amsmath,amssymb,bm}
\usepackage{mathtools}
\usepackage[colorlinks=true,linkcolor=blue,citecolor=blue,urlcolor=blue]{hyperref}

\newcommand{\dd}{\mathrm{d}}
\newcommand{\DD}{\mathrm{D}}
\newcommand{\gradq}{\nabla_{\!q}}
\newcommand{\gradx}{\nabla_{\!x}}
\newcommand{\lapq}{\nabla_{\!q}^{2}}
\newcommand{\lapx}{\nabla_{\!x}^{2}}
\newcommand{\PsiL}{\bm{\Psi}}
\newcommand{\xv}{\bm{x}}
\newcommand{\qv}{\bm{q}}
\newcommand{\kv}{\bm{k}}
\newcommand{\kone}{\bm{k}_1}
\newcommand{\ktwo}{\bm{k}_2}
\newcommand{\delD}{\delta_{\mathrm{D}}}
\newcommand{\tr}{\operatorname{tr}}

\title{Second-Order Lagrangian Perturbation Theory in an Einstein--de Sitter Universe}
\author{}
\date{}

\begin{document}
\maketitle

\section{Setup and conventions}

Work in comoving Eulerian coordinates $\xv$ and Lagrangian coordinates
$\qv$.  The trajectory of a cold dark matter fluid element is written
\begin{equation}
  \xv(\qv,t)=\qv+\PsiL(\qv,t),
\end{equation}
where $\PsiL$ is the Lagrangian displacement.  The peculiar velocity is
$\bm{u}=a\dot{\xv}=a\dot{\PsiL}$, with an overdot denoting a derivative
with respect to cosmic time $t$.

For an Einstein--de Sitter (EdS) universe,
\begin{equation}
  \Omega_m=1,\qquad \Omega_\Lambda=0,\qquad
  a(t)\propto t^{2/3},\qquad
  4\pi G\bar\rho = \frac{3}{2}H^2 ,
\end{equation}
where $H=\dot a/a$.  The comoving equation of motion and Poisson equation are
\begin{align}
  \ddot{\xv}+2H\dot{\xv}
    &= -\frac{1}{a^2}\gradx\Phi,\\
  \lapx\Phi
    &= 4\pi G a^2\bar\rho\,\delta
     = \frac{3}{2}a^2H^2\delta .
\end{align}
Thus the displacement acceleration
\begin{equation}
  \bm{A}\equiv \ddot{\PsiL}+2H\dot{\PsiL}
\end{equation}
satisfies
\begin{align}
  \gradx\cdot \bm{A} &= -\frac{3}{2}H^2\delta, \label{eq:euler-div}\\
  \gradx\times \bm{A} &= \bm{0}. \label{eq:euler-curl}
\end{align}
We assume the growing, irrotational solution and expand
\begin{equation}
  \PsiL = \PsiL^{(1)}+\PsiL^{(2)}+\cdots .
\end{equation}

\section{Mass conservation and the Jacobian}

Mass conservation gives
\begin{equation}
  \bar\rho\,[1+\delta(\xv,t)]\,\dd^3x
  = \bar\rho\,\dd^3q,
  \qquad
  1+\delta(\xv,t)=\frac{1}{J(\qv,t)},
\end{equation}
with
\begin{equation}
  J = \det\!\left(\frac{\partial x_i}{\partial q_j}\right)
    = \det\!\left(\delta_{ij}+\Psi_{i,j}\right),
  \qquad
  \Psi_{i,j}\equiv \frac{\partial \Psi_i}{\partial q_j}.
\end{equation}
Expanding the determinant in invariants of $\Psi_{i,j}$,
\begin{align}
  J &= 1+\mu_1+\mu_2+\mu_3,\\
  \mu_1 &= \Psi_{i,i},\\
  \mu_2 &= \frac{1}{2}\left[
      \left(\Psi_{i,i}\right)^2-\Psi_{i,j}\Psi_{j,i}
    \right],
\end{align}
where $\mu_3=\det(\Psi_{i,j})$ starts at cubic order.  To second order,
\begin{equation}
  \delta = \frac{1}{J}-1
  = -\mu_1-\mu_2+\mu_1^2+\mathcal{O}(\PsiL^3).
\end{equation}
Separating perturbative orders,
\begin{align}
  \delta^{(1)}
    &= -\Psi^{(1)}_{i,i}, \label{eq:delta-first}\\
  \delta^{(2)}
    &= -\Psi^{(2)}_{i,i}
       +\frac{1}{2}\left[
          \left(\Psi^{(1)}_{i,i}\right)^2
          +\Psi^{(1)}_{i,j}\Psi^{(1)}_{j,i}
        \right].
       \label{eq:delta-second}
\end{align}

\section{First order: the Zel'dovich displacement}

At first order $\gradx=\gradq+\mathcal{O}(\PsiL)$, so
Eq.~\eqref{eq:euler-div} gives
\begin{equation}
  \left(\frac{\dd^2}{\dd t^2}
       +2H\frac{\dd}{\dd t}
       -\frac{3}{2}H^2
  \right)
  \Psi^{(1)}_{i,i}=0.
  \label{eq:first-div-eq}
\end{equation}
The curl equation implies that the growing first-order displacement is
longitudinal.  We write
\begin{equation}
  \PsiL^{(1)}(\qv,t)=-D_1(t)\gradq\phi^{(1)}(\qv),
  \qquad
  \lapq\phi^{(1)}(\qv)=\delta_0(\qv),
  \label{eq:first-displacement}
\end{equation}
so that
\begin{equation}
  \delta^{(1)}(\qv,t)=D_1(t)\delta_0(\qv).
\end{equation}
The growth factor obeys
\begin{equation}
  \ddot D_1+2H\dot D_1-\frac{3}{2}H^2D_1=0.
  \label{eq:d1-eq}
\end{equation}
In EdS the two independent solutions are $D_1\propto a$ and
$D_1\propto a^{-3/2}$.  Keeping only the growing mode, we choose the
normalization
\begin{equation}
  D_1(t)=a(t).
\end{equation}

\section{Second-order equation}

To expand Eq.~\eqref{eq:euler-div} to second order, use
\begin{equation}
  \frac{\partial}{\partial x_i}
  =
  \left(\delta_{ji}-\Psi^{(1)}_{j,i}\right)
  \frac{\partial}{\partial q_j}
  +\mathcal{O}(\PsiL^2).
\end{equation}
Therefore
\begin{equation}
  \gradx\cdot\bm{A}
  =
  A^{(1)}_{i,i}
  +A^{(2)}_{i,i}
  -\Psi^{(1)}_{j,i}A^{(1)}_{i,j}
  +\mathcal{O}(\PsiL^3),
  \label{eq:lhs-second}
\end{equation}
where
\begin{equation}
  A^{(n)}_i
  =
  \ddot\Psi^{(n)}_i+2H\dot\Psi^{(n)}_i .
\end{equation}
Using Eq.~\eqref{eq:delta-second}, the second-order part of
Eq.~\eqref{eq:euler-div} is
\begin{multline}
  A^{(2)}_{i,i}-\Psi^{(1)}_{j,i}A^{(1)}_{i,j}
  =
  \frac{3}{2}H^2\Psi^{(2)}_{i,i} \\
  -\frac{3}{4}H^2
  \left[
      \left(\Psi^{(1)}_{i,i}\right)^2
      +\Psi^{(1)}_{i,j}\Psi^{(1)}_{j,i}
  \right].
  \label{eq:second-raw}
\end{multline}
For the growing first-order solution,
\begin{equation}
  A^{(1)}_{i,j}
  =
  \left(\ddot D_1+2H\dot D_1\right)
  \frac{\Psi^{(1)}_{i,j}}{D_1}
  =
  \frac{3}{2}H^2\Psi^{(1)}_{i,j},
\end{equation}
where Eq.~\eqref{eq:d1-eq} was used.  Substituting this into
Eq.~\eqref{eq:second-raw} yields
\begin{equation}
  \left(
    \frac{\dd^2}{\dd t^2}
    +2H\frac{\dd}{\dd t}
    -\frac{3}{2}H^2
  \right)
  \Psi^{(2)}_{i,i}
  =
  -\frac{3}{2}H^2\,
  \mu_2\!\left[\PsiL^{(1)}\right],
  \label{eq:second-div-eq}
\end{equation}
with
\begin{equation}
  \mu_2\!\left[\PsiL^{(1)}\right]
  =
  \frac{1}{2}
  \left[
    \left(\Psi^{(1)}_{i,i}\right)^2
    -\Psi^{(1)}_{i,j}\Psi^{(1)}_{j,i}
  \right].
  \label{eq:mu2-first}
\end{equation}

For irrotational growing initial conditions, the second-order transverse
source vanishes.  Hence the fastest-growing second-order displacement is also
longitudinal:
\begin{equation}
  \PsiL^{(2)}(\qv,t)=D_2(t)\gradq\phi^{(2)}(\qv).
  \label{eq:second-ansatz}
\end{equation}
Using Eq.~\eqref{eq:first-displacement},
\begin{equation}
  \mu_2\!\left[\PsiL^{(1)}\right]
  =
  D_1^2
  \sum_{i>j}
  \left[
    \phi^{(1)}_{,ii}\phi^{(1)}_{,jj}
    -\left(\phi^{(1)}_{,ij}\right)^2
  \right].
\end{equation}
Define the second-order potential by
\begin{equation}
  \lapq\phi^{(2)}
  =
  \sum_{i>j}
  \left[
    \phi^{(1)}_{,ii}\phi^{(1)}_{,jj}
    -\left(\phi^{(1)}_{,ij}\right)^2
  \right].
  \label{eq:phi2-poisson}
\end{equation}
Then Eq.~\eqref{eq:second-div-eq} reduces to the time equation
\begin{equation}
  \ddot D_2+2H\dot D_2-\frac{3}{2}H^2D_2
  =
  -\frac{3}{2}H^2D_1^2 .
  \label{eq:d2-eq}
\end{equation}

\section{EdS growing solution}

In EdS, set $D_1=a$.  Since
\begin{equation}
  \ddot a = -\frac{1}{2}H^2a,
\end{equation}
one finds
\begin{equation}
  \left(
    \frac{\dd^2}{\dd t^2}
    +2H\frac{\dd}{\dd t}
    -\frac{3}{2}H^2
  \right)a^2
  =
  \frac{7}{2}H^2a^2.
\end{equation}
The particular growing solution of Eq.~\eqref{eq:d2-eq} is therefore
\begin{equation}
  D_2(t)=-\frac{3}{7}D_1^2(t)
        =-\frac{3}{7}a^2(t).
  \label{eq:d2-eds}
\end{equation}
The homogeneous solutions are the same as at first order and are usually
discarded when constructing the fastest-growing 2LPT solution.

Combining the first- and second-order pieces, the standard EdS 2LPT map is
\begin{equation}
  \boxed{
  \xv(\qv,t)
  =
  \qv
  -D_1(t)\gradq\phi^{(1)}(\qv)
  +D_2(t)\gradq\phi^{(2)}(\qv)
  }
\end{equation}
with
\begin{align}
  \lapq\phi^{(1)} &= \delta_0,\\
  \lapq\phi^{(2)}
  &=
  \sum_{i>j}
  \left[
    \phi^{(1)}_{,ii}\phi^{(1)}_{,jj}
    -\left(\phi^{(1)}_{,ij}\right)^2
  \right],\\
  D_1 &= a,\qquad D_2=-\frac{3}{7}a^2.
\end{align}
Equivalently,
\begin{equation}
  \PsiL^{(2)}
  =
  -\frac{3}{7}a^2\gradq\phi^{(2)}.
\end{equation}

\section{Fourier-space form}

Use the Fourier convention
\begin{equation}
  f(\qv)
  =
  \int\frac{\dd^3k}{(2\pi)^3}\,
  e^{i\kv\cdot\qv}\,\tilde f(\kv).
\end{equation}
At a fixed time, absorb the powers of $D_1$ into the linear density
$\tilde\delta_L(\kv,t)=D_1(t)\tilde\delta_0(\kv)$.  Then
\begin{align}
  \tilde\Psi^{(1)}_i(\kv)
  &=
  i\,\frac{k_i}{k^2}\tilde\delta_L(\kv),\\
  \tilde\Psi^{(2)}_i(\kv)
  &=
  \frac{i}{2}
  \int\frac{\dd^3k_1}{(2\pi)^3}
       \frac{\dd^3k_2}{(2\pi)^3}
  (2\pi)^3\delD(\kv-\kone-\ktwo)\,
  L^{(2)}_i(\kone,\ktwo)\,
  \tilde\delta_L(\kone)\tilde\delta_L(\ktwo),
\end{align}
where the EdS second-order Lagrangian kernel is
\begin{equation}
  \boxed{
  \bm{L}^{(2)}(\kone,\ktwo)
  =
  \frac{3}{7}\,
  \frac{\kv}{k^2}
  \left[1-\mu_{12}^2\right]
  },
  \qquad
  \kv=\kone+\ktwo,\quad
  \mu_{12}=\frac{\kone\cdot\ktwo}{k_1k_2}.
\end{equation}

The corresponding second-order Eulerian density kernel obtained from the
Jacobian is the usual EdS result
\begin{equation}
  F_2(\kone,\ktwo)
  =
  \frac{5}{7}
  +\frac{1}{2}\mu_{12}
    \left(\frac{k_1}{k_2}+\frac{k_2}{k_1}\right)
  +\frac{2}{7}\mu_{12}^2.
\end{equation}
This agreement provides a useful check on the normalization and sign of
$D_2$.

\section{Planar fields with only \texorpdfstring{$k_z=0$}{kz = 0} modes}

The derivation above does not require statistical isotropy of the initial
density field.  Consider a three-dimensional EdS universe in which the density
field is independent of $z$, so that only transverse Fourier modes are present:
\begin{equation}
  \tilde\delta(\kv)
  \propto \delD(k_z).
\end{equation}
Equivalently, up to Fourier-convention factors and the sign convention in the
Dirac delta enforcing translation invariance, the two-point function may be
written schematically as
\begin{equation}
  \left\langle
    \tilde\delta(\kv)\tilde\delta(\kv')
  \right\rangle
  =
  P_{\mathrm{2D}}(k_\perp)\,
  \delD(k_z)\,\delD(k'_z)\,
  \delD^{(2)}(\kv_\perp+\kv'_\perp),
\end{equation}
where $\kv_\perp=(k_x,k_y)$ and $k_\perp=|\kv_\perp|$.

Because convolutions of $k_z=0$ modes remain at $k_z=0$, the LPT evolution
closes on the transverse plane.  The growing displacement has
\begin{equation}
  \Psi_z=0,\qquad
  \PsiL_\perp=\PsiL_\perp(\qv_\perp,t),
\end{equation}
and all spatial operators in the displacement potentials may be replaced by
their transverse versions,
\begin{equation}
  \gradq \rightarrow \nabla_{\perp},
  \qquad
  \lapq \rightarrow \nabla_{\perp}^{2}.
\end{equation}
The EdS time-dependent growth factors are unchanged, since their equations
come from the background expansion and the three-dimensional Poisson
normalization:
\begin{equation}
  D_1=a,\qquad D_2=-\frac{3}{7}a^2 .
\end{equation}
Thus the planar 2LPT map is
\begin{equation}
  \boxed{
  \xv_\perp(\qv_\perp,t)
  =
  \qv_\perp
  -a(t)\nabla_{\perp}\phi^{(1)}(\qv_\perp)
  -\frac{3}{7}a^2(t)\nabla_{\perp}\phi^{(2)}(\qv_\perp),
  \qquad
  z=q_z .
  }
\end{equation}
The potentials obey
\begin{align}
  \nabla_{\perp}^{2}\phi^{(1)}
    &= \delta_0(\qv_\perp),\\
  \nabla_{\perp}^{2}\phi^{(2)}
    &=
    \phi^{(1)}_{,xx}\phi^{(1)}_{,yy}
    -
    \left(\phi^{(1)}_{,xy}\right)^2 .
\end{align}
This is just the restriction of the three-dimensional 2LPT source
\begin{equation}
  \sum_{i>j}
  \left[
    \phi^{(1)}_{,ii}\phi^{(1)}_{,jj}
    -\left(\phi^{(1)}_{,ij}\right)^2
  \right]
\end{equation}
to fields for which all $z$ derivatives vanish.

In Fourier space the kernel is correspondingly restricted to transverse
vectors:
\begin{equation}
  \bm{L}^{(2)}_\perp(\kone,\ktwo)
  =
  \frac{3}{7}\,
  \frac{\kv_\perp}{k_\perp^2}
  \left[1-\mu_{12,\perp}^2\right],
  \qquad
  L_z^{(2)}=0,
\end{equation}
with
\begin{equation}
  \kv_\perp=\bm{k}_{1\perp}+\bm{k}_{2\perp},
  \qquad
  \mu_{12,\perp}
  =
  \frac{\bm{k}_{1\perp}\cdot\bm{k}_{2\perp}}
       {k_{1\perp}k_{2\perp}}.
\end{equation}

The important caveat is that this applies to a genuinely $z$-independent
three-dimensional density field.  A generic line-of-sight projection of a
three-dimensional density field is not an autonomous two-dimensional
gravitating system and cannot, in general, be evolved by simply replacing the
three-dimensional Laplacian with a two-dimensional one.  The normalization of
$P_{\mathrm{2D}}$ also depends on the Fourier and box-depth convention, but
that affects the statistical amplitude rather than the EdS growth factors.

\section*{References}

\begin{thebibliography}{9}
\bibitem{Zeldovich1970}
Ya.~B. Zel'dovich,
\emph{Gravitational instability: an approximate theory for large density perturbations},
Astronomy and Astrophysics \textbf{5}, 84 (1970).

\bibitem{Buchert1992}
T. Buchert,
\emph{Lagrangian perturbation theory: a key-model for large-scale structure},
Monthly Notices of the Royal Astronomical Society \textbf{254}, 729 (1992).

\bibitem{Bouchet1995}
F. R. Bouchet, S. Colombi, E. Hivon, and R. Juszkiewicz,
\emph{Perturbative Lagrangian approach to gravitational instability},
Astronomy and Astrophysics \textbf{296}, 575 (1995).

\bibitem{Bernardeau2002}
F. Bernardeau, S. Colombi, E. Gaztanaga, and R. Scoccimarro,
\emph{Large-scale structure of the Universe and cosmological perturbation theory},
Physics Reports \textbf{367}, 1 (2002).
\end{thebibliography}

\end{document}
